\chapter{結果の分析と考察}

\section{工学的作業量が形成した能力}

\figfull[.98\textwidth]{engineering_workload_dashboard.pdf}{初年度に確認できた作業量の全体像。データ資産、本番学習データ、学習探索、本番学習、運用モデル、視覚監査、強化学習研究の7群を示す。単位が異なるため合算はしない。}

これらは互いに孤立した数字ではなく、一つの工学連鎖を表す。専用ソフトで双腕教示を安定保存し、データ編集センターで軌跡を確認・整理・確定版として公開し、学習基盤で計算資源と設定を反復検討した。その後、本番運用モデルを現場へ投入し、アテンション記録で3視野への重み配分を監査し、別系統では強化学習の学習・保存・検証閉ループまで構築した。したがって初年度の価値は一度の展示に限られず、次年度の追加収集、比較評価、モデル改善へ再利用できる生産基盤にある。

\section{形成された主要能力}

\subsection{データから現場までの閉ループ}

51件の有効プロジェクトと2,413軌跡を管理するデータ基盤、25件の確定版データからなる本番学習、追跡可能な学習済みモデル、現場運用までを接続した。\ev{E02--E09} 新しいボトル型や難条件を追加する際、収集からやり直すのではなく、既存基盤へ追加して比較できる。

\subsection{多条件・多段階の教示資産}

単体・複数ボトル、方向変化、キャップなし、水入り、初期位置復帰、反転、自由回転キャップを含む。重点条件には反復サンプリングを適用した。この設計により、単一の定型把持ではなく、現場変化を含む多段階タスクの学習材料を整備した。

\subsection{双腕モデルと現場統合}

本番運用モデルは上部俯瞰、左右手首近景、双腕状態から将来50制御周期分の動作列を生成し、動作実行中に次の動作列を準備する。現場で全工程の反復を観測し、8,223件のアテンション記録で3視野への重み配分を監査できることを確認した。\ev{E06,E07,E10,E12}

\subsection{模倣学習から能力最適化への展開}

基礎モデルに加え、835軌跡、238,816学習区間、連続5ラウンド、固定オフライン検証からなる強化学習研究経路を形成した。まだ基礎モデルを置き換える段階ではないが、「実データ、人手評価、動作最適化、価値推定、保存、検証」という次年度の比較研究基盤がある。

\section{難しさ・投入作業・段階能力の対応}

\begin{table}[H]
\centering
\caption{タスク難度と投入作業の対応}
\small
\begin{tabularx}{\textwidth}{>{\bfseries}p{25mm}p{38mm}p{36mm}X}
\toprule
難しさ & 実施した作業 & 確認できる規模 & 形成した段階能力\\
\midrule
ランダムなボトル群と遮蔽 & 複数条件・多視点の収集と編集 & 51プロジェクト、2,413軌跡 & 複数ボトル場面で対象へ接近する基盤\\
口部と向きの判断 & 方向、キャップなし、反転、手首近景を重点収集 & 方向404、キャップなし158、反転127軌跡 & 複数配置条件への初期適応基盤\\
双腕接触協調 & 把持、保持、探索、回転、分別まで一貫教示 & 本番学習1,051軌跡、状態・動作14次元 & 全工程の双腕反復を現場で観測\\
多段階の連続制御 & 将来50制御周期分の動作列、先行計算、現場統合 & 50 Hz、ステップ19,000運用モデル & 単発動作を超える連続反復\\
大規模学習安定性 & バッチ・資源・設定を反復探索 & 41試行、14設定、6,059学習記録 & 再実行可能な学習・保存経路\\
視覚重みの監査 & 3視野のアテンション記録 & 9件の実行記録一覧、8,223サンプル & 視点別の重み配分を確認可能\\
次段階の最適化 & 強化学習データ、保存、固定検証 & 835軌跡、238,816学習区間、5ラウンド & 次年度のモデル改善比較基盤\\
\bottomrule
\end{tabularx}
\end{table}

各投入は単なる件数増加ではなく、特定の難所に対応し、次工程で再利用できる能力または資産を形成している。

\section{成果の成熟度}

\begin{table}[H]
\centering
\caption{主要成果の成熟度と次の定量化}
\begin{tabularx}{\textwidth}{>{\bfseries}p{48mm}p{31mm}X}
\toprule
成果 & 現在の成熟度 & 次の定量化\\
\midrule
データ・学習・保存・運用の閉ループ & 再確認可能 & 固定版と自動監査を継続\\
双腕動作生成と現場統合 & 再確認可能 & 入力取得から制御出力までの遅延と工程時間を記録\\
全タスクサイクル & 現場で観測 & 1本ごとの識別と工程別判定\\
1時間超、約2本/分 & 現場概算 & 開始終了、正味稼働時間、逐次計数\\
多条件適応・初期一般化 & データと現場兆候あり & 凍結した条件集合で反復試験\\
強化学習による動作最適化 & オフライン閉ループ成立 & 基礎モデルとの同条件実機比較\\
\bottomrule
\end{tabularx}
\end{table}

未測定項目を明確にすることは、構築済みの工学成果を弱めるものではない。全工程を実行するための統合機能とデータ生産基盤は構築済みであり、次年度は現場観測を反復可能な標準指標へ変換する段階である。

\section{重点的に改善する二つの現象}

\begin{table}[H]
\centering
\caption{現場現象、既存基盤、強化策}
\begin{tabularx}{\textwidth}{>{\bfseries}p{31mm}p{45mm}p{45mm}X}
\toprule
現象 & 既存基盤 & 次の強化 & 受入れ方法\\
\midrule
キャップなしでも開栓動作へ進む & キャップなし158軌跡を保有し、指示条件不足を特定 & キャップ有無以外をそろえた対照軌跡と条件付き指示を追加 & 開栓工程を正しく省略した割合を記録\\
一部を逆向きに把持 & 方向404軌跡、反転127軌跡を保有 & 口部方向、把持端、持上げ結果を付与 & 固定姿勢集合で正しい把持端を逐次集計\\
接触工程のばらつき & 手首視点に大きなアテンション重みが割り当てられる傾向があり、開栓教示を保有 & 接近、挟持、回転、離脱を分けて記録 & 工程別完了率、脱落、再試行を集計\\
\bottomrule
\end{tabularx}
\end{table}

両現象は既存基盤から改善可能な具体的入口が見つかっている。キャップの有無とボトル口の向きを対照条件で評価すれば、データ追加の効果を曖昧にせず確認できる。

\section{強化学習と高精度挿入への意味}

強化学習は、基礎模倣モデルを安定させた上で、同一条件の比較により増分を測るべきである。現在のオフライン約6.7\%改善は候補選別の根拠になるが、実機優位性の結論ではない。

空ボトルを水管へ挿入する次年度課題では、口部中心、管中心、軸方向、接触時の柔軟性が必要になる。現時点の「ミリメートル級」は目標精度であり、測定済み精度ではない。既存の多視点データ、双腕制御、強化学習、デジタルツイン研究を、幾何校正と段階別安全判定へ接続する必要がある。
